\documentclass[a4paper,12pt]{article}
\usepackage[utf8]{inputenc}
\usepackage{csquotes}
\usepackage[T1]{fontenc}
\usepackage[polish]{babel}
\usepackage{xcolor}
\usepackage{graphicx}
\usepackage{amsmath}
\usepackage{amssymb}
\usepackage{float}
\usepackage{listings}
% \usepackage[backend=biber, style=numeric]{biblatex} % Biber might require more setup
\usepackage[style=numeric]{biblatex} % Using bibtex backend by default if biber is tricky
\usepackage{hyperref}
\usepackage{enumitem}
\usepackage{longtable}
\usepackage{array}
\usepackage{booktabs}
\usepackage{tikz}

\usetikzlibrary{positioning, shapes.geometric, arrows.meta}

% \addbibresource{bibliografia.bib} % Assuming you have a bibliografia.bib file

\graphicspath{{images/}} % Make sure you have an 'images' folder

\lstset{
  basicstyle=\ttfamily\small,
  columns=flexible,
  keepspaces=true,
  showstringspaces=false,
  escapeinside={(*@}{@*)},
  literate={ą}{{\k{a}}}1
           {ć}{{\'{c}}}1
           {ę}{{\k{e}}}1
           {ł}{{\l{}}}1
           {ń}{{\'{n}}}1
           {ó}{{\'{o}}}1
           {ś}{{\'{s}}}1
           {ź}{{\'{z}}}1
           {ż}{{\.{z}}}1
           {Ą}{{\k{A}}}1
           {Ć}{{\'{C}}}1
           {Ę}{{\k{E}}}1
           {Ł}{{\L{}}}1
           {Ń}{{\'{N}}}1
           {Ó}{{\'{O}}}1
           {Ś}{{\'{S}}}1
           {Ź}{{\'{Z}}}1
           {Ż}{{\.{Z}}}1
           {"}{{\textquotedbl}}1
           {'}{{\textquotesingle}}1
           {`}{{\textasciigrave}}1
           {~}{{\textasciitilde}}1
           {^}{{\textasciicircum}}1
           {_}{{\textunderscore}}1
           {|}{{\textbar}}1
           {\{}{{\textbraceleft}}1
           {\}}{{\textbraceright}}1
           {[}{{[}}1
           {]}{{]}}1,
  language=SQL,
  showspaces=false,
  numbers=left,
  numberstyle=\tiny\color{gray},
  commentstyle=\color{green!60!black},
  keywordstyle=\color{blue},
  stringstyle=\color{red!80!black},
  breaklines=true,
  frame=tb, % Top and bottom frame for listings
  captionpos=b,
  tabsize=2
}

\hypersetup{
    colorlinks=true,
    linkcolor=blue,
    filecolor=magenta,
    urlcolor=cyan,
    pdftitle={Projekt hurtowni danych F1 - ETL},
    pdfauthor={Mikołaj Kubś},
    pdfpagemode=FullScreen,
}

\title{Projekt hurtowni danych do analizy wyników i czynników wpływających na osiągnięcia kierowców Formuły 1 w latach 1950-2023 \\ \vspace{0.5em} \large Etap II: Proces ETL}
\author{Mikołaj Kubś, 272662}
\date{26 maja 2025} % Or \today

\begin{document}

\maketitle
\tableofcontents
\newpage

\section{Wprowadzenie}
Niniejszy dokument stanowi raport z realizacji drugiego etapu projektu hurtowni danych, koncentrującego się na procesie ETL (Extract, Transform, Load). Celem tego etapu było zasilenie zdefiniowanych w Etapie I tabel wymiarów i faktów danymi pochodzącymi z plików CSV, dotyczących wyników i czynników wpływających na osiągnięcia kierowców Formuły 1 w latach 1950-2023. Proces został zaimplementowany przy użyciu narzędzia SQL Server Integration Services (SSIS).

\section{Architektura i Narzędzia Procesu ETL}

\subsection{Struktura Procesu ETL}
Proces ETL został zorganizowany w ramach jednego projektu SSIS, zawierającego dedykowane pakiety dla ładowania poszczególnych tabel wymiarów oraz tabeli faktów. Dodatkowo, zaimplementowano pakiet główny (master package) koordynujący uruchamianie poszczególnych pakietów w odpowiedniej kolejności (najpierw wymiary, następnie fakty), aby zapewnić integralność referencyjną.

Struktura pakietów SSIS (przykładowa):
\begin{itemize}
  \item \texttt{Load\_Dim\_Circuit.dtsx}
  \item \texttt{Load\_Dim\_Driver.dtsx}
  \item \texttt{Load\_Dim\_Constructor.dtsx}
  \item \texttt{Load\_Dim\_Race.dtsx}
  \item \texttt{Load\_Dim\_Time.dtsx}
  \item \texttt{Load\_Dim\_Status.dtsx}
  \item \texttt{Load\_Dim\_Weather.dtsx}
  \item \texttt{Load\_Fact\_Result.dtsx}
  \item \texttt{Master\_F1\_DWH\_Load.dtsx}
\end{itemize}

\subsection{Mapa Logiczna ETL (Uproszczona)}
Poniżej przedstawiono uproszczoną mapę logiczną ilustrującą przepływ danych od plików źródłowych, przez transformacje, do tabel docelowych w hurtowni.
\begin{figure}[H]
  \centering
  \begin{tikzpicture}[
      node distance=1.5cm and 2.5cm,
      source/.style={rectangle, draw, fill=blue!20, minimum height=1cm, minimum width=2.5cm, text centered},
      transform/.style={ellipse, draw, fill=orange!20, minimum height=1cm, text centered},
      dwh_table/.style={cylinder, shape border rotate=90, draw, fill=green!20, aspect=0.5, minimum height=1.2cm, minimum width=1cm, text centered},
      arrow/.style={-Latex}
    ]
    % Sources
    \node[source] (csv_circuits) {circuits.csv};
    \node[source, right=of csv_circuits] (csv_races) {races.csv};
    \node[source, right=of csv_races] (csv_results) {results.csv};
    \node[source, right=of csv_results] (csv_weather) {weather.csv};

    % Staging/Transform (Simplified)
    \node[transform, below=of csv_circuits] (trans_circuits) {DFT: Dim\_Circuit};
    \node[transform, below=of csv_races] (trans_races) {DFT: Dim\_Race};
    \node[transform, below=of csv_results] (trans_results) {DFT: Fact\_Result};
    \node[transform, below=of csv_weather] (trans_weather) {DFT: Dim\_Weather};

    % DWH Tables
    \node[dwh_table, below=of trans_circuits] (dim_circuit) {Dim\_Circuit};
    \node[dwh_table, right=of dim_circuit] (dim_race) {Dim\_Race};
    \node[dwh_table, right=of dim_race] (fact_result) {Fact\_Result};
    \node[dwh_table, right=of fact_result] (dim_weather) {Dim\_Weather};

    % Arrows
    \draw[arrow] (csv_circuits) -- (trans_circuits);
    \draw[arrow] (csv_races) -- (trans_races);
    \draw[arrow] (csv_results) -- (trans_results);
    \draw[arrow] (csv_weather) -- (trans_weather);

    \draw[arrow] (trans_circuits) -- (dim_circuit);
    \draw[arrow] (trans_races) -- (dim_race);
    \draw[arrow] (trans_results) -- (fact_result);
    \draw[arrow] (trans_weather) -- (dim_weather);

    % Connections to Fact_Result (example)
    \draw[arrow] (dim_circuit.north) .. controls +(north:1cm) and +(north:1cm) .. (fact_result.west);
    \draw[arrow] (dim_race.north) -- (fact_result.north west);
    \draw[arrow] (dim_weather.north) .. controls +(north:1cm) and +(north:1cm) .. (fact_result.east);

  \end{tikzpicture} % <<--- THIS IS THE CRUCIAL FIX
  \caption{Uproszczona mapa logiczna procesu ETL.}
  \label{fig:etl_map}
\end{figure}
*(Uwaga: Pełna mapa logiczna powinna obejmować wszystkie źródła, kluczowe transformacje i tabele docelowe, w tym lookupy i agregacje.)*

\subsection{Narzędzia}
Głównym narzędziem wykorzystanym do implementacji procesów ETL był SQL Server Integration Services (SSIS) wchodzący w skład Microsoft SQL Server. Pakiety SSIS zostały utworzone przy użyciu SQL Server Data Tools (SSDT) for Visual Studio.

\section{Implementacja Procesów ETL}

\subsection{Ekstrakcja Danych (Extract)}
Dane źródłowe w formacie CSV były wczytywane przy użyciu komponentu \texttt{Flat File Source} w SSIS. Kluczowe aspekty konfiguracji tego etapu to:
\begin{itemize}
  \item Ustawienie strony kodowej na \texttt{65001 (UTF-8)} dla wszystkich plików źródłowych, aby zapewnić poprawne odczytywanie znaków specjalnych.
  \item Poprawna identyfikacja separatorów kolumn, kwalifikatorów tekstu oraz obsługa wierszy nagłówkowych.
  \item Wstępne mapowanie typów danych na poziomie źródła, z uwzględnieniem późniejszych konwersji.
\end{itemize}

\subsection{Transformacja Danych (Transform)}
Etap transformacji obejmował szereg operacji mających na celu oczyszczenie, integrację i przygotowanie danych do załadowania do tabel docelowych.
\begin{enumerate}
  \item \textbf{Konwersja Typów Danych:} Komponent \texttt{Data Conversion} był intensywnie wykorzystywany do jawnej konwersji typów danych z plików źródłowych (często tekstowych) na typy docelowe zdefiniowane w modelu hurtowni (np. \texttt{INT}, \texttt{DATE}, \texttt{DECIMAL}, \texttt{NVARCHAR} lub \texttt{VARCHAR} z odpowiednią stroną kodową). Szczególną uwagę zwrócono na konwersję:
        \begin{itemize}
          \item Kolumn tekstowych (string) na \texttt{Unicode string [DT\_WSTR]} w SSIS, co odpowiada typowi \texttt{NVARCHAR} w SQL Server, w celu uniknięcia problemów z różnymi stronami kodowymi i zapewnienia obsługi znaków międzynarodowych. Alternatywnie, dla kolumn docelowych \texttt{VARCHAR}, konwersja na \texttt{string [DT\_STR]} z docelową stroną kodową (np. 1250).
          \item Reprezentacji dat i czasów na odpowiednie typy (\texttt{DT\_DATE}, \texttt{DT\_DBTIME2}).
          \item Wartości numerycznych.
        \end{itemize}
  \item \textbf{Obsługa Wartości NULL:} Wartości reprezentujące NULL w plikach CSV (np. `\\N`) były identyfikowane i konwertowane na rzeczywiste wartości \texttt{NULL} w przepływie danych przy użyciu komponentu \texttt{Derived Column}.
  \item \textbf{Lookup'y Kluczy Zastępczych:} Podczas ładowania tabeli faktów (\texttt{Fact\_Result}), komponent \texttt{Lookup Transformation} był używany do pobierania kluczy zastępczych (surrogate keys) z wcześniej załadowanych tabel wymiarów, bazując na kluczach naturalnych (np. \texttt{driverId}, \texttt{raceId}).
  \item \textbf{Obliczenia Pochodne:} Komponent \texttt{Derived Column} służył do tworzenia nowych kolumn na podstawie istniejących, np.:
        \begin{itemize}
          \item Obliczenie wieku kierowcy w momencie wyścigu (\texttt{AgeAtRace}).
          \item Obliczenie zmiany pozycji (\texttt{PositionOrderChange}).
        \end{itemize}
  \item \textbf{Agregacja Danych:}
        \begin{itemize}
          \item \textbf{Dane pogodowe (\texttt{Dim\_Weather}):} Minutowe dane pogodowe z pliku \texttt{F1 Weather (2023-2018).csv} były agregowane do poziomu pojedynczego wyścigu (po \texttt{Year} i \texttt{RoundNumber}) przy użyciu komponentu \texttt{Aggregate}. Obliczano wartości średnie, minimalne, maksymalne oraz sumy (np. dla opadów).
          \item \textbf{Dane o pit-stopach:} Informacje o liczbie i łącznym czasie pit-stopów dla każdego kierowcy w wyścigu były agregowane z pliku \texttt{pit\_stops.csv} i dołączane podczas ładowania tabeli faktów (np. poprzez wcześniejszą agregację do tabeli stagingowej i lookup).
        \end{itemize}
  \item \textbf{Kategoryzacja Danych Pogodowych:} Po agregacji, zagregowane wartości liczbowe dla danych pogodowych (np. temperatura, wilgotność, prędkość wiatru) były przekształcane na predefiniowane kategorie tekstowe (np. "Zimno", "Umiarkowanie", "Gorąco") przy użyciu wyrażeń warunkowych w komponencie \texttt{Derived Column}, zgodnie z założeniami projektu. Te kategorie zostały następnie załadowane do \texttt{Dim\_Weather}.
\end{enumerate}

\subsection{Ładowanie Danych (Load)}
Przetransformowane dane były ładowane do docelowych tabel w hurtowni SQL Server przy użyciu komponentu \texttt{OLE DB Destination}.
\begin{itemize}
  \item Wykorzystano tryb szybkiego ładowania (\texttt{Table or view - fast load}).
  \item Zapewniono poprawne mapowanie kolumn z przepływu danych SSIS do kolumn tabel docelowych.
  \item Kolejność ładowania (najpierw wymiary, potem fakty) była zarządzana przez pakiet główny.
\end{itemize}

\subsection{Wyzwania i Rozwiązania}
Podczas implementacji napotkano typowe wyzwania związane z procesami ETL:
\begin{itemize}
  \item \textbf{Niezgodność Stron Kodowych:} Początkowe problemy z konfliktem stron kodowych (np. 65001 z Flat File Source vs 1250 domyślne dla OLE DB Destination) zostały rozwiązane poprzez:
        \begin{enumerate}
          \item Konsekwentne ustawienie kodowania \texttt{65001 (UTF-8)} dla komponentów \texttt{Flat File Source}.
          \item Zastosowanie komponentu \texttt{Data Conversion} do jawnej konwersji kolumn tekstowych na typ \texttt{Unicode string [DT\_WSTR]} (dla docelowych kolumn \texttt{NVARCHAR}) lub \texttt{string [DT\_STR]} z określoną docelową stroną kodową (dla kolumn \texttt{VARCHAR}).
        \end{enumerate}
  \item \textbf{Synchronizacja Metadanych:} Zmiany w strukturze przepływu danych (np. po dodaniu \texttt{Data Conversion}) wymagały odświeżenia metadanych w komponencie \texttt{OLE DB Destination} i ponownego zmapowania kolumn.
  \item \textbf{Błędy Konwersji Danych w Źródle:} Błędy typu "Text was truncated or one or more characters had no match in the target code page" pojawiające się w \texttt{Flat File Source} wskazywały na niepoprawną konfigurację strony kodowej dla pliku źródłowego. Problem rozwiązano poprzez ustawienie strony kodowej na \texttt{65001 (UTF-8)}, co odpowiadało rzeczywistemu kodowaniu plików.
\end{itemize}

\subsection{Przyrostowe Zasilanie (Planowane)}
Dla tabeli faktów \texttt{Fact\_Result} zaplanowano mechanizm przyrostowego zasilania. Polega on na przechowywaniu identyfikatora ostatnio załadowanego rekordu (np. \texttt{resultId} w tabeli kontrolnej) i ładowaniu w kolejnych uruchomieniach ETL tylko nowszych danych.
*(Uwaga: Opisać, czy zostało to w pełni zaimplementowane, czy jest w planach na dalszy rozwój.)*

\section{Podsumowanie i Wnioski z Etapu ETL}
Proces ETL został pomyślnie zaimplementowany, umożliwiając zasilenie hurtowni danych Formuły 1 przetworzonymi i zintegrowanymi danymi. Kluczowe było właściwe zarządzanie typami danych i stronami kodowymi oraz logiczne rozplanowanie transformacji, w tym agregacji i kategoryzacji danych pogodowych. Zastosowanie SSIS pozwoliło na zbudowanie modularnego i zarządzalnego przepływu danych. Rozwiązane problemy z konwersją danych podkreślają znaczenie dokładnej analizy danych źródłowych i właściwej konfiguracji narzędzi ETL.

Kolejnym etapem będzie budowa kostki OLAP na bazie przygotowanej hurtowni danych oraz przeprowadzenie analiz wielowymiarowych.

\newpage
% \section*{Dodatek A: Przykładowe fragmenty kodu / konfiguracji SSIS}
% (Opcjonalnie: można tu dodać zrzuty ekranu konfiguracji kluczowych komponentów SSIS, np. Data Conversion dla pogody, Lookup, fragment mapy logicznej z SSIS)


% \printbibliography % If using biblatex

\end{document}